\documentclass{article}
\usepackage[spanish]{babel}
\usepackage{multicol}
\usepackage{amsmath}
\usepackage{amsfonts}
\usepackage{enumerate}
\usepackage[utf8]{inputenc}
\usepackage{amssymb}
\usepackage{graphicx}
\usepackage{wasysym}
\usepackage{moodle}
\usepackage[usenames]{color}
\newcommand{\radio}{\ooalign{\hidewidth$\bullet$\hidewidth\cr$\ocircle$}}
\newcommand{\R}{$\mathbb{R}^3$ }
\newcommand{\wi}{\displaystyle \int}
\newtheorem{theorem}{Theorem}
\usepackage[left=2.00cm, right=2.00cm, top=2.00cm, bottom=2.00cm]{geometry}
\title{Parcial 2}
\author{Departamento de matemáticas}
\date{\today}

\begin{document}
	\maketitle
	\begin{quiz}{}
	\begin{cloze}{}
		Considere la integral:
		\[\displaystyle\int_{-\infty}^{\infty}f(x)dx\]
		donde $f(x) = \displaystyle\frac{e^{-x^2}\sin(x)}{x\sqrt{x^2+1}}$
		Dicha integral se debe de escribir de la forma:
		\[\displaystyle\int_{-\infty}^{-1}f(x)dx + \displaystyle\int_{A}^{B}f(x)dx+\displaystyle\int_{B}^{C}f(x)dx+ \displaystyle\int_{1}^{\infty}f(x)dx\]
	    \begin{numerical}[points=5]{}
	    	el valor A es \item -1
	    \end{numerical}
         \begin{numerical}[points=5]{}
        	el valor B es \item 0
        \end{numerical}
         \begin{numerical}[points=5]{}
        	el valor C es \item 1
        \end{numerical}
		\begin{numerical}[points=18]{}
			Realizando la sustituci\'on apropiada y usando una cuadratura Gaussiana de 3 nodos la integral con tres cifras decimales da aproximadamente \item[tolerance=0.003] 1.422
		\end{numerical}
	\end{cloze}
	\begin{cloze}{}
	Considere la integral:
	\[\displaystyle\int_{-\infty}^{\infty}f(x)dx\]
	donde $f(x) = \displaystyle\frac{2e^{-(x-1)^2}\sin(x-1)}{(x-1)\sqrt{(x-1)^2+1}}$
	Dicha integral se debe de escribir de la forma:
	\[\displaystyle\int_{-\infty}^{-1}f(x)dx + \displaystyle\int_{-1}^{A}f(x)dx +  \displaystyle\int_{A}^{B}f(x)dx + \displaystyle\int_{2}^{\infty}f(x)dx\]
	\begin{numerical}[points=8]{}
		el valor A es \item 1
	\end{numerical}
	\begin{numerical}[points=7]{}
		el valor B es \item 2
	\end{numerical}
	\begin{numerical}[points=18]{}
		Realizando la sustituci\'on apropiada y usando una cuadratura Gaussiana de 3 nodos la integral con tres cifras decimales da aproximadamente \item[tolerance=0.003] 2.866
	\end{numerical}
\end{cloze}
\begin{cloze}{}
	Considere la integral:
	\[\displaystyle\int_{-\infty}^{\infty}f(x)dx\]
	donde $f(x) = \displaystyle\frac{3e^{-(x+1)^2}\sin(x+1)}{(x+1)\sqrt{(x+1)^2+1}}$
	Dicha integral se debe de escribir de la forma:
	\[\displaystyle\int_{-\infty}^{-2}f(x)dx + \displaystyle\int_{A}^{B}f(x)dx+\displaystyle\int_{B}^{1}f(x)dx+ \displaystyle\int_{1}^{\infty}f(x)dx\]
	\begin{numerical}[points=8]{}
		el valor A es \item -2
	\end{numerical}
	\begin{numerical}[points=7]{}
		el valor B es \item -1
	\end{numerical}
	\begin{numerical}[points=18]{}
		Realizando la sustituci\'on apropiada y usando una cuadratura Gaussiana de 3 nodos la integral con tres cifras decimales da aproximadamente \item[tolerance=0.003] 4.299
	\end{numerical}
\end{cloze}
\begin{cloze}{}
	Considere la integral:
	\[\displaystyle\int_{-\infty}^{\infty}f(x)dx\]
	donde $f(x) = \displaystyle\frac{2e^{-(x+2)^2}\sin(x+2)}{(x+2)\sqrt{(x+2)^2+1}}$
	Dicha integral se debe de escribir de la forma:
	\[\displaystyle\int_{-\infty}^{-3}f(x)dx + \displaystyle\int_{A}^{B}f(x)dx+\displaystyle\int_{B}^{C}f(x)dx+ \displaystyle\int_{-1}^{1}f(x)dx+ \displaystyle\int_{1}^{\infty}f(x)dx\]
	\begin{numerical}[points=5]{}
		el valor A es \item -3
	\end{numerical}
	\begin{numerical}[points=5]{}
		el valor B es \item -2
	\end{numerical}
	\begin{numerical}[points=5]{}
		el valor C es \item -1
	\end{numerical}
	\begin{numerical}[points=18]{}
		Realizando la sustituci\'on apropiada y usando una cuadratura Gaussiana de 3 nodos la integral con tres cifras decimales da aproximadamente \item[tolerance=0.003] 2.864
	\end{numerical}
\end{cloze}
\begin{cloze}{}
	Considere la integral:
	\[\displaystyle\int_{-\infty}^{\infty}f(x)dx\]
	donde $f(x) = \displaystyle\frac{4e^{-(x-0.5)^2}\sin(x-0.5)}{(x-0.5)\sqrt{(x-0.5)^2+1}}$
	Dicha integral se debe de escribir de la forma:
	\[\displaystyle\int_{-\infty}^{-1}f(x)dx + \displaystyle\int_{-1}^{A}f(x)dx +  \displaystyle\int_{A}^{B}f(x)dx + \displaystyle\int_{1}^{\infty}f(x)dx\]
	\begin{numerical}[points=8]{}
		el valor A es \item 0.5
	\end{numerical}
	\begin{numerical}[points=7]{}
		el valor B es \item 1
	\end{numerical}
	\begin{numerical}[points=18]{}
		Realizando la sustituci\'on apropiada y usando una cuadratura Gaussiana de 3 nodos la integral con tres cifras decimales da aproximadamente \item[tolerance=0.003] 5.641
	\end{numerical}
\end{cloze}
    \end{quiz}
\end{document}